\documentclass{article}
\usepackage[utf8]{inputenc}
\usepackage[T1]{fontenc}
\usepackage{geometry}
\geometry{
    a4paper,
    left=2.5cm,
    right=2.5cm,
    top=2.5cm,
    bottom=2.5cm
}
\usepackage{graphicx}
\usepackage{float}
\usepackage{hyperref}
\usepackage[english]{babel}
\usepackage{amsmath}
\usepackage{amssymb}
\usepackage{array}
\usepackage{booktabs}
\usepackage{caption}
\usepackage{subcaption}
\usepackage{enumitem}
\usepackage{fancyhdr}
\usepackage{titlesec}
\usepackage{url}
\usepackage{textcomp}
\usepackage{xcolor}
\usepackage{listings}
\usepackage{multirow}
\usepackage{longtable}
\usepackage{tabularx}
\usepackage{times}

% Set up page style
\pagestyle{fancy}
\fancyhf{}
\rfoot{\thepage}
\renewcommand{\headrulewidth}{0pt}

% Title formatting
\titleformat{\chapter}[display]{\normalfont\bfseries}{\chaptertitle\thechapter}{0pt}{\Huge}
\titlespacing*{\chapter}{0pt}{50pt}{40pt}

% Document info
\title{Template Document}
\author{Academic Support Agent}
\date{\today}

\begin{document}
\maketitle

\maketitle

%%% define code
\defverbatim[colored]\lstI{
	\begin{lstlisting}[language=C++,basicstyle=\ttfamily,keywordstyle=\color{red}]
	int main() {
	// Define variables at the beginning
	// of the block, as in C:
	CStash intStash, stringStash;
	int i;
	char* cp;
	ifstream in;
	string line;
	[...]
	\end{lstlisting}
}
%%%%%%%%%%%%%%%%%%%%%%%%%%%%%%%%
% ----------- FRAME ------------
%%%%%%%%%%%%%%%%%%%%%%%%%%%%%%%%

\section{Basics}
\subsection{Blocks}
\begin{frame}[c]{Blocks}
	
The blocks are shown below
 block
\begin{block}{Regular Block}
	Content of a regular block
\end{block}

\begin{exampleblock}{Example Block}
	Content of an example block
\end{exampleblock}

\begin{alertblock}{Alert block}
	Content of an alert block
\end{alertblock}

\end{frame}	

\subsection{Enumerate \& Overlays}

\begin{frame}[c]{Enumerate \& Overlays}
	
{\large An Example of \texttt{enumerate}}
	\begin{enumerate}[<+->]
		\item First item
		\item Second item
		\item Third item
	\end{enumerate}
\vfill
{\large An Example of \texttt{itemize}}	
	\begin{itemize}[<+->]
	\item First item
	\item Second item
	\item Third item
	\end{itemize}
\end{frame}	

\subsection{Two columns}
\begin{frame}[c]{Two columns}
	
\begin{columns}[onlytextwidth]
	\column{0.4\textwidth}
	Content for column one
	\begin{equation}
	E = mc^2
	\end{equation}
	\column{0.4\textwidth}	
	Content for column two
	\begin{equation}
	F=ma
	\end{equation}
\end{columns}
\end{frame}	

\begin{frame}{ }
\begin{columns}
	\column{0.5\textwidth}
		\begin{figure}
		\centering
		% Requires \usepackage{graphicx}
		\includegraphics[width=5cm]{logo}
		\caption{ }
	\end{figure}
	\column{0.55\textwidth}
	\small
	 
\end{columns}
\end{frame}


\subsection{Figures}
\begin{frame}[c]{Figures}
	\begin{figure}
		\centering
		\begin{tikzpicture}
		\draw [help lines,very thick] (0,0) grid (5,4);
		\end{tikzpicture}
		\caption{Credits to Ti\textit{k}Z}
	\end{figure}
\end{frame}	

\begin{frame}{Figures}
 :
\begin{figure}[h]
    \centering
    \subfigure[ 1]{
    \includegraphics[scale=0.022]{logo.png}
    }
    \subfigure[ 2]{
    \includegraphics[scale=0.022]{logo.png}
    }
    \subfigure[ 3]{
    \includegraphics[scale=0.022]{logo.png}
    }
    \subfigure[ 4]{
    \includegraphics[scale=0.022]{logo.png}
    }
    \caption{demo}
\end{figure}
\end{frame}

\subsection{Code Demo}
\begin{frame}{A Listings Demo}{C++}
	\lstI
\end{frame}

\section{References}

\begin{frame}{BibTex }
     \cite{9492070} 
\end{frame}

\begin{frame}{References}
     % 



\end{frame}
% Thank you page
\beamertemplateshadingbackground{structure.fg!90}{structure.fg}
\begin{frame}[plain]
	\vfill
	\centering
	{
		\centering \Huge \color{white} Thank you for your attention!\\[10pt]Questions?
	}
	\vfill
\end{frame}
\end{document}